\chapter{Drei Beispiele zur Reichweite}

Die folgenden endlichen Tafeln erläutern Voraussetzungen und Grenzen des
Hauptsatzes. Ihre Rechnungen verwenden keine weiteren Strukturresultate.
Verschiedene Buchstaben in den jeweils angegebenen Trägern bezeichnen
verschiedene Elemente; insbesondere sind alle als Nichtnullwerte bezeichneten
Elemente verschieden von \(0\).

\section{Eine Dreierhalbgruppe mit zwei Produktwerten}
\hypertarget{nullproducts.example.three}{}

Sei \(S=\{0,c,a\}\). Wir definieren eine binäre Operation durch
\[
\begin{array}{c|ccc}
 \cdot&0&c&a\\\hline
 0&0&0&0\\
 c&0&0&0\\
 a&0&0&c
\end{array}
\]
Jede Tabellenstelle besitzt genau einen Wert in \(S\); damit ist die
Operation eine Funktion \(S\times S\to S\). Für beliebige
\(x,y,z\in S\) gilt \(xy,yz\in\{0,c\}\). Die ersten beiden
Zeilen und Spalten ergeben deshalb vollständig
\[
\begin{array}{c|cc}
 &xy=0&xy=c\\\hline
 (xy)z&0z=0&cz=0
\end{array}
\qquad
\begin{array}{c|cc}
 &yz=0&yz=c\\\hline
 x(yz)&x0=0&xc=0
\end{array}
\]
Die Fälle sind jeweils erschöpfend. Also ist \((xy)z=x(yz)\), und
\(S\) ist eine Halbgruppe. Zeile und Spalte von \(0\) zeigen, dass
\(0\) ihre Null ist. Jedes Produkt liegt in \(\{0,c\}\); wegen
\(00=0\) und \(aa=c\) treten beide Werte auf. Damit ist
\(S^2=\{0,c\}\).

Für nichtleere Teilmengen \(A,B\subseteq S\) gilt:
\[
 AB=\{0\}\quad\Longleftrightarrow\quad
 a\notin A\ \lor\ a\notin B.
\]
Enthalten beide Mengen \(a\), so liegt \(aa=c\ne0\) im Produkt,
das daher nicht \(\{0\}\) sein kann. Fehlt \(a\) in einer der
Mengen, ist jedes Faktorenpaar verschieden von \((a,a)\); die Tafel
gibt für jedes solche Paar den Wert \(0\). Wegen der Nichtleerheit der
beiden Mengen tritt \(0\) tatsächlich auf, also ist das Produkt
\(\{0\}\). Damit sind beide Richtungen gezeigt.

Die beiden Profile sind folglich
\[
\begin{aligned}
 p&=\{\{0\},\{c\},\{0,c\}\},\\
 q&=\{\{a\},\{0,a\},\{c,a\},\{0,c,a\}\}.
\end{aligned}
\]
Jede Menge des Profils \(p\) besteht alle Nulltests; eine Menge des
Profils \(q\) besteht genau die Tests mit einer Menge ohne \(a\).
Die beiden Profile sind verschieden, wie der Test mit \(\{a\}\)
zeigt. Es gilt \(p<q\), und die Mengen aus dem Zählargument sind
\[
\begin{array}{c|c|c|c|c}
 \text{Profil}&L&D&C&|L|\\\hline
 p&\PowNE{\{0,c\}}&\{0,c\}&\{0,c\}&3\\
 q&\PowNE S&S&\{a\}&7
\end{array}
\]
So gewinnt \(|L_p|=2^{|D_p|}-1\) zuerst \(|D_p|=2\) und
\(|D_q|=3\) zurück. Die disjunkte Zerlegung
\(D_q=C_p\mathbin{\dot\cup}C_q\) liefert dann \(|C_q|=1\).
Hier liegen die beiden ausgezeichneten Elemente \(0,c\) im selben
Profil: Der Zweipunktfall im Hauptbeweis ist also bereits für dieses
kleine Beispiel notwendig.

\section{Ein Potenzhalbgruppenautomorphismus kann Einermengen vertauschen}
\hypertarget{nullproducts.example.singletons}{}

Sei \(N=\{0,a\}\), und sei jedes Produkt in \(N\) gleich \(0\).
Die Operation ist abgeschlossen, und für jedes Tripel sind beide
Klammerungen gleich \(0\). Sie definiert somit eine Halbgruppe mit Null.
Auf ihren drei nichtleeren Teilmengen definieren wir
\[
 \sigma(\{0\})=\{0\},\qquad
 \sigma(\{a\})=\{0,a\},\qquad
 \sigma(\{0,a\})=\{a\}.
\]
Die dreimalige Auswertung ergibt für jede der drei Mengen
\(\sigma(\sigma(A))=A\); \(\sigma\) ist also seine eigene inverse
Abbildung und damit bijektiv. Für beliebige nichtleere \(A,B\subseteq N\)
ist jedes einzelne Produkt \(ab\) gleich \(0\), und mindestens ein
solches Faktorenpaar existiert. Daher ist \(AB=\{0\}\). Auch
\(\sigma(A)\) und \(\sigma(B)\) sind nichtleer, weshalb
\[
 \sigma(AB)=\sigma(\{0\})=\{0\}=\sigma(A)\sigma(B).
\]
Damit ist \(\sigma\) ein Automorphismus der Potenzhalbgruppe. Die
Einermenge \(\{a\}\) wird dabei auf die Zweiermenge \(\{0,a\}\)
abgebildet. Das Beispiel besitzt nur einen Produktwert und widerlegt keine
Voraussetzung oder Folgerung des Hauptsatzes. Es widerlegt die für dessen
Beweis unnötige Annahme, jeder Potenzhalbgruppenisomorphismus müsse
Einermengen erhalten.

\section{Drei Produktwerte: gleiche Nullrelation, verschiedene Halbgruppen}
\hypertarget{nullproducts.example.threevalues}{}

Auf \(E=\{0,a,b,u,v\}\) seien die beiden Operationen vollständig durch
die folgenden Tafeln definiert:
\[
\begin{array}{c|ccccc}
 \cdot&0&a&b&u&v\\\hline
 0&0&0&0&0&0\\
 a&0&u&v&0&0\\
 b&0&v&u&0&0\\
 u&0&0&0&0&0\\
 v&0&0&0&0&0
\end{array}
\qquad
\begin{array}{c|ccccc}
 \ast&0&a&b&u&v\\\hline
 0&0&0&0&0&0\\
 a&0&u&u&0&0\\
 b&0&v&u&0&0\\
 u&0&0&0&0&0\\
 v&0&0&0&0&0
\end{array}
\]
Jede Stelle ist mit genau einem Element von \(E\) besetzt, also sind
beide Operationen wohldefinierte Funktionen \(E\times E\to E\).
Schreiben wir \(\circ\) für eine der beiden Operationen und
\(Z=\{0,u,v\}\), so zeigen alle Stellen der betreffenden Tafel
\(x\circ y\in Z\). Die drei Zeilen und Spalten zu \(Z\) sind
vollständig null; für jedes \(z\in Z\) und jedes \(t\in E\) gilt
\(z\circ t=0=t\circ z\). Für beliebige \(x,y,t\in E\) folgt daher
\[
 (x\circ y)\circ t=0=x\circ(y\circ t).
\]
Dies deckt alle Tripel ab und beweist die Assoziativität beider
Operationen. Die Zeile und Spalte zu \(0\) beweisen die Nulleigenschaft.
Die Werte \(0,u,v\) treten bei \(0\circ0\), \(a\circ a\) und
\(b\circ a\) auf; weitere Produktwerte gibt es laut Tafel nicht.
Beide Produktquadrate sind also \(Z\).

Die Nullrelation ist für beide Operationen genau dieselbe:
\[
 x\circ y=0\quad\Longleftrightarrow\quad
 (x,y)\notin\{a,b\}\times\{a,b\}.
\]
Denn im ausgezeichneten Viererblock stehen ausschließlich die von \(0\)
verschiedenen Werte \(u,v\), außerhalb stehen ausschließlich Nullen.
Damit stimmen auch alle Nulltests nichtleerer Teilmengen und somit ihre
Nullproduktprofile überein.

Die Halbgruppe \((E,\cdot)\) ist kommutativ: Außerhalb des Viererblocks
bleiben die Nullen bei Vertauschung der Faktoren erhalten; innerhalb
stimmen die beiden Diagonalprodukte ohnehin mit sich selbst überein und
die beiden anderen Stellen haben beide den Wert \(v\). Hingegen ist
\((E,\ast)\) nicht kommutativ, weil \(a\ast b=u\ne v=b\ast a\).

Ein Isomorphismus erhält die Kommutativität: Für beliebige Zielelemente
wähle man Urbilder \(x,y\); dann gilt
\(f(x)f(y)=f(xy)=f(yx)=f(y)f(x)\). Die beiden Halbgruppen sind
somit nicht isomorph. Ebenso sind ihre Potenzhalbgruppen nicht isomorph:
Links kommutieren alle Einzelprodukte und damit sämtliche Mengenprodukte;
rechts kommutieren schon \(\{a\},\{b\}\) nicht. Das Beispiel zeigt,
dass gleiche Nullproduktrelationen bei drei Produktwerten die volle
Multiplikation nicht festlegen; eine Isomorphie der Potenzhalbgruppen
wird nicht behauptet.
